The original frozen campaign recorded 72/72 joint results accepted by its declared oracle, corresponding to 216/216 family acceptance checks. Within each treatment, each class has 24 accepted outputs across six projects and four cycles. No retries were needed on these supported inputs; four separately counted refusal controls rejected cyclic dependencies, missing evidence, private data and job-count overflow. There were 24 observed parent invocations and 12 child invocations. These counters establish use, not economic benefit.

The stable repeat reduced six-project CPU from \textbf{2.271 seconds} for full rebuild to \textbf{0.845 seconds} for the flywheel, while conventional caching used \textbf{0.822 seconds}. On the two fresh-change cycles, flywheel and conventional caching each rebuilt only \textbf{18 of 132 source cards} and reused 114. Their total CPU was \textbf{4.069} and \textbf{4.033 seconds}, respectively. The flywheel was approximately \textbf{0.9\% slower} on this recorded fresh slice. Six local instances do not establish a population timing difference; they establish no demonstrated superiority over the strong incumbent.

\begin{figure}[H]\centering\includegraphics[width=\linewidth]{figures/production.pdf}
\caption{Recorded native research. CPU includes the joint local pipeline, native/portfolio planning, checks and persisted artifacts; it is not complete Agent cost. The source-card reduction is also achieved by ordinary caching. Each point summarizes six controlled projects.}\end{figure}

\begin{table}[H]\small\centering\begin{tabular}{lrrr}\toprule Cycle & Full rebuild CPU (s) & Flywheel CPU (s) & Cache CPU (s)\\\midrule
Initial & 2.337 & 2.451 & 2.489 \\
Stable repeat & 2.271 & 0.845 & 0.822 \\
Source/duration change & 2.401 & 2.369 & 2.273 \\
Dependency/calendar repair & 1.674 & 1.700 & 1.760 \\
\bottomrule\end{tabular}\caption{Observed joint CPU; one cost belongs to three deliverables and is not counted three times.}\end{table}

Observed median local wall times for a joint flywheel request were 361.9, 122.0, 360.1, 230.0 milliseconds for the initial, stable, source-change and dependency-change cycles, respectively. They exclude unmetered authoring, human review, customer communication and the unavailable live executor path. No active human-minute measurement is available.

Child descriptor construction across six bindings consumed 0.421 seconds wall time and 0.466 seconds CPU in addition to service execution. These measurements cover authored-algorithm binding and checking, not the human/model effort that designed or corrected the algorithm. Original initialization, developmental failures and refused inputs are retained separately.

\subsection{Post-freeze defects and the corrective edition}
Independent-role adversarial review after the original freeze found stale path-only metadata, cross-project predecessor reuse, inadequately checked construction origin and inconsistent deletion handling. The original oracle also failed some tampered metadata/digest probes. Those controls were outside the original fixed-inventory confirmation inputs. The original acceptance totals are preserved with that scope; they do not qualify arbitrary changed projects. A separately frozen corrective adapter and fresh edge evaluation are reported below. No historical result was overwritten or promoted to the new source identity.

The corrective v1.1 adapter uses path- and provenance-aware content keys, rejects cross-project predecessors, checks immutable construction origin and preserves an accepted child registry for subsequent use. After developmental review, its own freeze bound 36 files at 14:19:54 UTC. A fresh qualification then produced \textbf{25 accepted outputs and three expected refusals}: \textbf{28/28 cases passed}, with 3,265 strengthened Python assertions. A separately authored saved-output audit passed 3,885 assertions. These finite edge cases qualify the correction only within their scope; they do not replace the original 72-run timing campaign or establish general correctness. The browser demo uses v1.1. The economic measurement bridge retains the original scoped timing results without claiming corrected timing superiority.


The new change-impact and cross-artifact certificate is a real implemented interface. The strong conventional alternatives can produce the same deliverables under the declared bounds. Consequently, \textbf{previously impractical work becoming feasible under the same budgets is not established}. Discovery-process improvement and self-acceleration are also not established.
